\section{Fibered categories}
\source{fga-explained:page-0051}\source{fga-explained:page-0052}\source{fga-explained:page-0053}\source{fga-explained:page-0054}\source{fga-explained:page-0055}\source{fga-explained:page-0056}\source{fga-explained:page-0057}\source{fga-explained:page-0058}\source{fga-explained:page-0059}\source{fga-explained:page-0060}\source{fga-explained:page-0061}\source{fga-explained:page-0062}\source{fga-explained:page-0063}\source{fga-explained:page-0064}\source{fga-explained:page-0065}\source{fga-explained:page-0066}\source{fga-explained:page-0067}\source{fga-explained:page-0068}\source{fga-explained:page-0069}\source{fga-explained:page-0070}\source{fga-explained:page-0071}\source{fga-explained:page-0072}\source{fga-explained:page-0073}\source{fga-explained:page-0074}\source{fga-explained:page-0075}\source{fga-explained:page-0076}

\begin{definition}[fibered category]\label{fga-def-3-1}\dcref{3.1}\source{fga-explained:page-0051}\uses{}
A category $\mathcal X$ fibered over $\mathcal C$ is equipped with a functor
$p:\mathcal X\to\mathcal C$ such that for every arrow $f:U\to V$ and object
$\eta$ over $V$ there is a cartesian arrow $f^*\eta\to\eta$ over $f$.
\end{definition}\begin{definition}[cartesian arrow]\label{fga-def-3-5}\dcref{3.5}\source{fga-explained:page-0052}\uses{fga-def-3-1}
An arrow $\alpha:\xi\to\eta$ over $f:U\to V$ is cartesian if every arrow to
$\eta$ over a composite $fg$ factors uniquely through $\alpha$ over $g$.
\end{definition}
\begin{proposition}\label{fga-prop-3-7}\dcref{3.7}\source{fga-explained:page-0053}\uses{fga-def-3-5}
Cartesian arrows are stable under composition and contain all identities; an
arrow isomorphism over an identity is cartesian.
\end{proposition}

\begin{definition}[cleavage]\label{fga-def-3-8}\dcref{3.8}\source{fga-explained:page-0053}\uses{fga-def-3-5}
A cleavage chooses, for every $f$ and $\eta$, one cartesian pullback arrow
$f^*\eta\to\eta$.  A cleavage is normalized when it chooses identities as
identities and split when $(gf)^*=f^*g^*$ on the nose.
\end{definition}
\begin{definition}[fiber and base change]\label{fga-def-3-9}\dcref{3.9}\source{fga-explained:page-0054}\uses{fga-def-3-8}
The fiber $\mathcal X(U)$ consists of objects over $U$ and arrows over
$\operatorname{id}_U$.  A cleavage induces pullback functors $f^*:
\mathcal X(V)\to\mathcal X(U)$.
\end{definition}
\begin{proposition}\label{fga-prop-3-11}\dcref{3.11}\source{fga-explained:page-0055}\uses{fga-def-3-5,fga-def-3-9}
For a fibered category, cartesian pullbacks are unique up to a unique vertical
isomorphism, and a chosen cleavage gives coherent canonical comparison
isomorphisms $(gf)^*\simeq f^*g^*$.
\end{proposition}
\begin{proof}
Apply the universal property twice to the two cartesian arrows.  The resulting
vertical maps are inverse by uniqueness, and applying the same argument to a
triple composite proves the cocycle coherence.
\end{proof}

\begin{definition}[pseudofunctor]\label{fga-def-3-12}\dcref{3.12}\source{fga-explained:page-0055}\uses{fga-def-3-9}
A pseudofunctor $\mathcal C^{\mathrm{op}}\to\mathbf{Cat}$ assigns categories
$\mathcal X(U)$ and pullback functors, together with coherent isomorphisms
$f^*g^*\simeq(gf)^*$ and $\operatorname{id}^*\simeq\operatorname{id}$.
\end{definition}
\begin{proposition}\label{fga-prop-3-13}\dcref{3.13}\source{fga-explained:page-0056}\uses{fga-def-3-12,fga-prop-3-11}
Fibered categories with cleavage and pseudofunctors are equivalent descriptions;
the Grothendieck construction sends a pseudofunctor to a fibered category.
\end{proposition}

\begin{definition}[stable class]\label{fga-def-3-16}\dcref{3.16}\source{fga-explained:page-0060}\uses{}
A class $\mathcal P$ of arrows is stable if it is invariant under isomorphism
and base change, with the relevant fiber products existing.
\end{definition}\begin{definition}[fibered in groupoids]\label{fga-def-3-21}\dcref{3.21}\source{fga-explained:page-0063}\uses{fga-def-3-1}
$\mathcal X\to\mathcal C$ is fibered in groupoids when every vertical arrow is
invertible.
\end{definition}
\begin{proposition}\label{fga-prop-3-22}\dcref{3.22}\source{fga-explained:page-0063}\uses{fga-def-3-21,fga-def-3-5}
For a fibered category in groupoids, an arrow is cartesian exactly when it is
universal among arrows with its base map; cartesian arrows are closed under
vertical isomorphism and composition.
\end{proposition}
\begin{corollary}\label{fga-cor-3-23}\dcref{3.23}\source{fga-explained:page-0063}\uses{fga-prop-3-22,fga-prop-3-13}
A fibered category in groupoids is determined up to equivalence by its fibers and
pullback functors with their coherence isomorphisms.
\end{corollary}

\begin{definition}[fibered in sets]\label{fga-def-3-24}\dcref{3.24}\source{fga-explained:page-0064}\uses{fga-def-3-1}
A fibered category in sets has discrete fibers: every vertical arrow is an
identity.  It is therefore equivalent to a contravariant set-valued functor.
\end{definition}
\begin{proposition}\label{fga-prop-3-25}\dcref{3.25}\source{fga-explained:page-0064}\uses{fga-def-3-24}
The category of elements of a contravariant set-valued functor is fibered in
sets; conversely every fibered category in sets arises this way.
\end{proposition}
\begin{proposition}\label{fga-prop-3-26}\dcref{3.26}\source{fga-explained:page-0065}\uses{fga-def-3-5,fga-def-3-21}
A morphism of fibered categories is cartesian when it preserves cartesian arrows;
equivalences can be tested fiberwise for a fibered category in groupoids.
\end{proposition}

\begin{definition}[fibered category associated with an object]\label{fga-def-3-29}\dcref{3.29}\source{fga-explained:page-0068}\uses{fga-def-3-1}
For $X\in\mathcal C$, the fibered category $\mathcal C/X\to\mathcal C$ has
objects arrows $U\to X$ and arrows commutative triangles; cartesian arrows are
cartesian squares in $\mathcal C$.
\end{definition}
\begin{definition}[2-Yoneda map]\label{fga-def-3-31}\dcref{3.31}\source{fga-explained:page-0069}\uses{fga-def-3-29}
For a fibered category $\mathcal X$, a morphism $\mathcal C/X\to\mathcal X$
over $\mathcal C$ is a family of objects and compatible pullback arrows; this is
the fibered analogue of a universal element.
\end{definition}
\begin{proposition}\label{fga-prop-3-32}\dcref{3.32}\source{fga-explained:page-0069}\uses{fga-def-3-31}
Morphisms $\mathcal C/X\to\mathcal X$ over $\mathcal C$ are equivalent to
objects of the fiber $\mathcal X(X)$.
\end{proposition}
\begin{definition}[functor of arrows]\label{fga-def-3-33}\dcref{3.33}\source{fga-explained:page-0070}\uses{fga-def-3-1}
The functor of arrows of a fibered category records, over $U$, arrows in the
fiber over $U$ and their pullbacks; it is itself fibered over $\mathcal C$.
\end{definition}
\begin{proposition}\label{fga-prop-3-35}\dcref{3.35}\source{fga-explained:page-0071}\uses{fga-def-3-33}
The arrow construction commutes with base change and its cartesian arrows are
exactly those whose underlying source and target arrows are cartesian.
\end{proposition}
\begin{theorem}\label{fga-thm-3-45}\dcref{3.45}\source{fga-explained:page-0074}\uses{fga-prop-3-32,fga-def-3-33}
For a fibered category $\mathcal X\to\mathcal C$ there is a canonical fibered
category of morphisms into $\mathcal X$; the 2-Yoneda correspondence identifies
its objects over $U$ with objects of $\mathcal X(U)$ and its arrows with
compatible morphisms.
\end{theorem}
\begin{proof}
Construct the category by Grothendieck's category-of-elements procedure and
verify the cartesian lifting property using the universal property of chosen
pullbacks.  The correspondence is evaluation at the identity of the base.
\end{proof}

\begin{definition}[equivariant object]\label{fga-def-3-46}\dcref{3.46}\source{fga-explained:page-0075}\uses{fga-def-2-15,fga-def-3-1}
For a group object $G$ acting on the base, an object $\xi\in\mathcal X(X)$ is
$G$-equivariant when it is supplied with an isomorphism between its translates
along the action satisfying the identity and cocycle conditions.
\end{definition}
\begin{proposition}\label{fga-prop-3-47}\dcref{3.47}\source{fga-explained:page-0075}\uses{fga-def-3-46,fga-def-3-33}
An equivariant structure on $\xi$ is equivalent to a lift of the action map to
the category of arrows of $\mathcal X$.
\end{proposition}
\begin{proposition}\label{fga-prop-3-48}\dcref{3.48}\source{fga-explained:page-0076}\uses{fga-prop-3-47}
Equivariant structures pull back along equivariant maps and form a category;
their morphisms are exactly equivariant morphisms in the fiber.
\end{proposition}
